\subsubsection{Tempered Distributions}\label{subsec:tempered-distributions}

It is possible to define different spaces of test functions and, correspondingly, of distributions. In particular, for $\Omega=\mathbb{R}^m$, it is natural to replace the requirement of compact support by one of rapid decay at infinity. This leads to the following definition.

\begin{Definition}
\label{def:rapidly-decreasing-test-space}
Let $\mathcal{S}(\mathbb{R}^m)$ be the space of all complex-valued functions on $\mathbb{R}^m$ which are of class $C^\infty$ and such that $|\bm{x}|^k|D^\alpha\phi(\bm{x})|$ is bounded for every $k\in\mathbb{N}$ and every multi-index $\alpha$. We say that a sequence $\phi_n$ in $\mathcal{S}(\mathbb{R}^m)$ converges to $\phi$ if the derivatives of all orders of the $\phi_n$ converge uniformly to those of $\phi$ and the constants $C_{k\alpha}$ in the bounds $|\bm{x}|^k|D^\alpha\phi_n(\bm{x})|\leq C_{k\alpha}$ can be chosen independently of $n$.
\end{Definition}

Obviously, $\mathcal{D}(\mathbb{R}^m)$ is a subspace of $\mathcal{S}(\mathbb{R}^m)$. Moreover, $\mathcal{D}(\mathbb{R}^m)$ is dense in $\mathcal{S}(\mathbb{R}^m)$. To see this, let $e(\bm{x})$ be a $C^\infty$-function which is equal to 1 in the unit ball and vanishes outside the ball of radius 2. Let $e_n(\bm{x})=e(\bm{x}/n)$. Then, for any $f\in\mathcal{S}(\mathbb{R}^m)$, we have $f=\lim_{n\to\infty}fe_n$.

We now define the tempered distributions to be continuous linear functionals on $\mathcal{S}$.

\begin{Definition}
\label{def:tempered-distribution}
A \textbf{tempered distribution} on $\mathbb{R}^m$ is a linear mapping $\phi\mapsto(f,\phi)$ from $\mathcal{S}(\mathbb{R}^m)$ to $\mathbb{C}$ with the continuity property that $(f,\phi_n)\to(f,\phi)$ if $\phi_n\to\phi$ in $\mathcal{S}(\mathbb{R}^m)$. The set of all tempered distributions is denoted by $\mathcal{S}^{\prime}(\mathbb{R}^m)$. We say that $f_n\to f$ in $\mathcal{S}^{\prime}(\mathbb{R}^m)$ if $(f_n,\phi)\to(f,\phi)$ for every $\phi\in\mathcal{S}(\mathbb{R}^m)$.
\end{Definition}

Clearly, every tempered distribution defines a distribution by restriction. Moreover, if two tempered distributions agree as elements of $\mathcal{D}^{\prime}(\mathbb{R}^m)$, they also agree as elements of $\mathcal{S}^{\prime}(\mathbb{R}^m)$; this follows from the fact that $\mathcal{D}(\mathbb{R}^m)$ is dense in $\mathcal{S}(\mathbb{R}^m)$. Hence $\mathcal{S}^{\prime}(\mathbb{R}^m)$ is a linear subspace of $\mathcal{D}^{\prime}(\mathbb{R}^m)$. Moreover, convergence in $\mathcal{S}^{\prime}(\mathbb{R}^m)$ obviously implies convergence in $\mathcal{D}^{\prime}(\mathbb{R}^m)$.
